\documentclass{article}
\usepackage[utf8]{inputenc}
\usepackage{hyperref}
\begin{document}

\title{Hot Interconnects Special Issue of IEEE Micro: How to Submit and Additional Information}
\maketitle

\noindent If you have any questions that are not answered here please do not hesitate to reach out 
to the special issue editors for HotI: Whit Schonbein (\href{mailto:wwschon@sandia.gov}{wwschon@sandia.gov}) or 
Joseph Schuchart (\href{mailto:joseph.schuchart@stonybrook.edu}{joseph.schuchart@stonybrook.edu}).

\bigskip

\noindent To submit:
\begin{enumerate}
\item Go to: \url{https://www.computer.org/csdl/magazine/mi}
\item Select the "Write for Us $\rightarrow$ Submit a Manuscript" option. Note: 
The Hot Interconnects IEEE Special Issue will not appear under the "Call for Papers" as this is an invited 
special issue. 
\item Select the "Start Submission" button.
\item Select the "Article Type" button, and from the dropdown list, select "SI: Hot Interconnects". 
If the special issue is not an option, it is because those responsible for adding it to the website have not yet done so, 
so check back later.
\item After the previous step follow the instructions as requested by the IEEE Micro submission system.
\end{enumerate}

\bigskip

\noindent Additional information:
\begin{enumerate}
\item IEEE Micro is considered by IEEE to be a `magazine' rather than a journal. Consequently, while the 
articles remain technical, they are typically shorter than a journal article. The content should be written for a broader audience.
\item IEEE Micro enforces a word limit of 6000 words, where each figure counts as 
250 words, as well as imposing limits on the number of references and length of 
abstract. See \url{https://www.computer.org/csdl/magazine/mi/write-for-us/14289} for more information.
\item To distinguish an IEEE Micro article from the conference proceeding version, 
there should be `30\%' new content. Because the length of the IEEE Micro article is 
likely to be \emph{less} than the original conference paper, this means a submission may 
require significant editing to both trim the length while adding new content. 
\item To distinguish an IEEE Micro article from the conference proceeding version, please 
use a \emph{different title} for the IEEE Micro article.
\item It is permissible for the author list to change for the IEEE Micro submission 
(either through re-ordering or adding or subtracting authors) in comparison to the 
conference version.
\item Use of a specific IEEE Micro article LaTeX template is not required for submitting 
the manuscript for review. One can use, e.g., a generic LaTeX article class. If one wants 
to use the IEEE Micro template, it can be downloaded using the IEEE template selector 
located here: \url{https://journals.ieeeauthorcenter.ieee.org/create-your-ieee-journal-article/authoring-tools-and-templates/tools-for-ieee-authors/ieee-article-templates/}.
\item Author biographies are required, and follow a prescribed format (\url{https://www.computer.org/publications/author-resources#author-biographies})


\end{enumerate}

\end{document}
